\chapter{Weak Solutions, Part II}

\section{Local Boundedness}

\begin{theorem}\label{hanlin:weaksol-part2-local-boundedness}\dcref{Theorem 4.1}\group{ch04-weak-solutions-part-ii}
\source{han-lin-lecture-notes:page-0073,han-lin-lecture-notes:page-0074,han-lin-lecture-notes:page-0075,han-lin-lecture-notes:page-0076,han-lin-lecture-notes:page-0077,han-lin-lecture-notes:page-0078,han-lin-lecture-notes:page-0079,han-lin-lecture-notes:page-0080,han-lin-lecture-notes:page-0081}
Suppose $a_{ij}\in L^\infty(B_1)$ and $c\in L^q(B_1)$ for some
$q>n/2$, with
\[
a_{ij}(x)\xi_i\xi_j\geq\lambda|\xi|^2
\qquad(x\in B_1,\ \xi\in\R^n),
\]
and
\[
|a_{ij}|_{L^\infty}+\|c\|_{L^q}\leq\Lambda.
\]
Let $u\in H^1(B_1)$ be a subsolution in the sense that
\[
(*)\qquad
\int_{B_1}a_{ij}D_i uD_j\varphi+cu\varphi
\leq\int_{B_1}f\varphi
\]
for every nonnegative $\varphi\in H^1_0(B_1)$. If $f\in L^q(B_1)$,
then $u^+\in L^\infty_{\mathrm{loc}}(B_1)$ and, for every
$\theta\in(0,1)$ and $p>0$,
\[
\sup_{B_\theta}u^+
\leq C\left\{
\frac{1}{(1-\theta)^{n/p}}\|u^+\|_{L^p(B_1)}
+\|f\|_{L^q(B_1)}\right\},
\]
where $C$ depends only on $n,\lambda,\Lambda,p$, and $q$.
\end{theorem}
\begin{proof}
\sketch It suffices first to prove the estimate for $p=2$ and
$\theta=1/2$. Put $\bar u=u^++k$ and truncate it by
\[
\bar u_m=
\begin{cases}
\bar u,&u<m,\\
k+m,&u\geq m.
\end{cases}
\]
For $\beta\geq0$ and $0\leq\eta\in C^1_0(B_1)$, use
\[
\varphi=\eta^2(\bar u_m^\beta\bar u-k^{\beta+1})
\]
in $(*)$. With $c_0=|c|+|f|/k$ and
$w=\bar u_m^{\beta/2}\bar u$, the ellipticity and Cauchy inequalities give
\[
\int|D(w\eta)|^2
\leq C(1+\beta)\int w^2|D\eta|^2
+C(1+\beta)\int c_0w^2\eta^2.
\]
Choose $k=\|f\|_{L^q}$ when $f\not\equiv0$; otherwise take $k>0$ and
later let $k\downarrow0$. Then $\|c_0\|_{L^q}\leq\Lambda+1$.
Interpolation between $L^2$ and the Sobolev exponent absorbs the last term,
so, for some exponent $a=a(n,q)>0$,
\[
\left(\int|\eta w|^{2\chi}\right)^{1/\chi}
\leq C(1+\beta)^a
\int(|D\eta|^2+\eta^2)w^2,
\]
where $\chi=n/(n-2)>1$ if $n>2$, while any fixed $\chi>1$ may be used
if $n=2$. Letting $m\to\infty$, setting $\gamma=\beta+2$, and choosing a
cutoff between $B_r$ and $B_R$ gives the reverse H\"older estimate
\begin{equation}
\|\bar u\|_{L^{\gamma\chi}(B_r)}
\leq\left(C\frac{(\gamma-1)^a}{(R-r)^2}\right)^{1/\gamma}
\|\bar u\|_{L^\gamma(B_R)}.
\end{equation}
Iterate this reverse H\"older estimate with $\gamma_i=2\chi^i$ and
$r_i=1/2+2^{-i-1}$. Since $\sum i/\chi^i<\infty$,
\begin{equation}
\sup_{B_{1/2}}u^+
\leq C\bigl(\|u^+\|_{L^2(B_1)}+\|f\|_{L^q(B_1)}\bigr).
\end{equation}

For $0<R\leq1$, apply the preceding $L^2$ bound to
$\widetilde u(y)=u(Ry)$, whose coefficients and forcing are
\[
\widetilde a(y)=a(Ry),\qquad
\widetilde c(y)=R^2c(Ry),\qquad
\widetilde f(y)=R^2f(Ry).
\]
This yields, for $p\geq2$,
\begin{equation}
\sup_{B_{R/2}}u^+
\leq C\left\{R^{-n/p}\|u^+\|_{L^p(B_R)}
+R^{2-n/q}\|f\|_{L^q(B_R)}\right\}.
\end{equation}
Covering $B_{\theta R}$ by balls of radius $(1-\theta)R$ gives the stated
$\theta$ dependence. If $0<p<2$, interpolate
$\|u^+\|_{L^2(B_R)}$ between $L^p$ and $L^\infty$ in the scaled estimate
to obtain
\[
F(r)\leq\frac12F(R)
+\frac{C}{(R-r)^{n/p}}\|u^+\|_{L^p(B_1)}
+C\|f\|_{L^q(B_1)},
\]
where $F(t)=\|u^+\|_{L^\infty(B_t)}$.
Lemma~\ref{hanlin:weaksol-part2-iteration-lemma} removes the
$F(R)$ term.
\end{proof}

\begin{remark}\label{hanlin:weaksol-part2-bounded-test-function}\dcref{Remark 4.2}\group{ch04-weak-solutions-part-ii}
\source{han-lin-lecture-notes:page-0080}
\uses{hanlin:weaksol-part2-local-boundedness}
If $u$ is bounded, the truncation in the proof may be omitted and one may use
\[
\varphi=\eta^2(\bar u^{\beta+1}-k^{\beta+1})
\in H^1_0(B_1)
\]
for $\beta\geq0$ and $0\leq\eta\in C^1_0(B_1)$.
\end{remark}

\begin{lemma}\label{hanlin:weaksol-part2-iteration-lemma}\dcref{Lemma 4.3}\group{ch04-weak-solutions-part-ii}
\source{han-lin-lecture-notes:page-0081}
Let $f\geq0$ be bounded on $[\tau_0,\tau_1]$, with $\tau_0\geq0$, and
suppose
\[
f(t)\leq\theta f(s)+\frac{A}{(s-t)^\alpha}+B
\qquad(\tau_0\leq t<s\leq\tau_1),
\]
where $\theta\in[0,1)$. Then
\[
f(t)\leq c\left(\frac{A}{(s-t)^\alpha}+B\right)
\qquad(\tau_0\leq t<s\leq\tau_1),
\]
where $c$ depends only on $\alpha$ and $\theta$.
\end{lemma}
\begin{proof}
\sketch Fix $t<s$ and choose $\tau\in(0,1)$ with
$\theta\tau^{-\alpha}<1$. For
\[
t_0=t,\qquad
t_{i+1}=t_i+(1-\tau)\tau^i(s-t),
\]
iteration gives
\[
f(t)\leq\theta^kf(t_k)+
\left(\frac{A}{(1-\tau)^\alpha(s-t)^\alpha}+B\right)
\sum_{i=0}^{k-1}\theta^i\tau^{-i\alpha}.
\]
Let $k\to\infty$ and sum the geometric series.
\end{proof}

\begin{theorem}\label{hanlin:weaksol-part2-high-integrability}\dcref{Theorem 4.4}\group{ch04-weak-solutions-part-ii}
\source{han-lin-lecture-notes:page-0082,han-lin-lecture-notes:page-0083}
\uses{hanlin:weaksol-part2-local-boundedness}
Assume $n\geq3$. Suppose $a_{ij}\in L^\infty(B_1)$ and
$c\in L^{n/2}(B_1)$ satisfy
\[
\lambda|\xi|^2\leq a_{ij}(x)\xi_i\xi_j\leq\Lambda|\xi|^2
\qquad(x\in B_1,\ \xi\in\R^n).
\]
Let $u\in H^1(B_1)$ satisfy
\[
\int_{B_1}a_{ij}D_i uD_j\varphi+cu\varphi
\leq\int_{B_1}f\varphi
\]
for every nonnegative $\varphi\in H^1_0(B_1)$. If
$f\in L^q(B_1)$ for
$q\in[2n/(n+2),n/2)$, then
$u^+\in L^{q^*}_{\mathrm{loc}}(B_1)$, where
\[
\frac1{q^*}=\frac1q-\frac2n.
\]
Moreover,
\[
\|u^+\|_{L^{q^*}(B_1)}
\leq C\bigl(\|u^+\|_{L^2(B_1)}+\|f\|_{L^q(B_1)}\bigr),
\]
where $C$ depends only on $n,\lambda,\Lambda,q$, and
\[
\varepsilon(K)=
\left(\int_{\{|c|>K\}}|c|^{n/2}\right)^{2/n}.
\]
\end{theorem}
\begin{proof}
\sketch Put $\bar u=u^+$, truncate it by
$\bar u_m=\min\{\bar u,m\}$, and test with
$\eta^2\bar u^\beta\bar u_m$. With $\chi=n/(n-2)$, the energy and Sobolev
inequalities give
\[
\left(\int\eta^{2\chi}\bar u_m^{\beta\chi}\bar u^{2\chi}
\right)^{1/\chi}
\leq C(1+\beta)\left\{
\int(|D\eta|^2+|c|\eta^2)\bar u_m^\beta\bar u^2
+\int|f|\eta^2\bar u_m^\beta\bar u\right\}.
\]
Split the $c$ term at level $K$:
\[
\int|c|\eta^2\bar u_m^\beta\bar u^2
\leq K\int\eta^2\bar u_m^\beta\bar u^2
+\varepsilon(K)
\left(\int(\eta^2\bar u_m^\beta\bar u^2)^\chi\right)^{1/\chi}.
\]
For bounded $\beta$, choose $K$ so that the last term is absorbed. The
forcing term is absorbed by H\"older and Young inequalities provided
\[
\beta+2\leq\frac{q(n-2)}{n-2q}=\frac{q^*}{\chi}.
\]
Thus, for $2\leq\gamma=\beta+2\leq q^*/\chi$ and $0<r<R\leq1$,
\begin{equation}
\|\bar u\|_{L^{\chi\gamma}(B_r)}
\leq C\left((R-r)^{-2/\gamma}
\|\bar u\|_{L^\gamma(B_R)}+\|f\|_{L^q(B_1)}\right).
\end{equation}
Finitely iterate the preceding estimate through $2,2\chi,2\chi^2,\ldots$
until the exponent
reaches $q^*/\chi$, and apply the reverse H\"older estimate once more to
obtain the asserted
$L^{q^*}$ estimate.
\end{proof}

\section{H\"older Continuity}

Let
\[
Lu=-D_i(a_{ij}(x)D_j u)
\]
on $B_1$, where
\[
\lambda|\xi|^2\leq a_{ij}(x)\xi_i\xi_j\leq\Lambda|\xi|^2
\qquad(x\in B_1,\ \xi\in\R^n).
\]

\begin{definition}\label{hanlin:weaksol-part2-homogeneous-subsolution}\dcref{Definition 4.5}\group{ch04-weak-solutions-part-ii}
\source{han-lin-lecture-notes:page-0084}
A function $u\in H^1_{\mathrm{loc}}(B_1)$ is a subsolution, respectively a
supersolution, of $Lu=0$ if
\[
\int_{B_1}a_{ij}D_i uD_j\varphi\leq0,
\qquad\text{respectively }\geq0,
\]
for every nonnegative $\varphi\in H^1_0(B_1)$.
\end{definition}

\begin{lemma}\label{hanlin:weaksol-part2-convex-composition}\dcref{Lemma 4.6}\group{ch04-weak-solutions-part-ii}
\source{han-lin-lecture-notes:page-0084}
\uses{hanlin:weaksol-part2-homogeneous-subsolution}
Let $\Phi\in C^1_{\mathrm{loc}}(\R)$ be convex.
\begin{enumerate}
\item[(i)] If $u$ is a subsolution and $\Phi'\geq0$, then
$\Phi(u)$ is a subsolution whenever $\Phi(u)\in H^1_{\mathrm{loc}}(B_1)$.
\item[(ii)] If $u$ is a supersolution and $\Phi'\leq0$, then
$\Phi(u)$ is a subsolution whenever $\Phi(u)\in H^1_{\mathrm{loc}}(B_1)$.
\end{enumerate}
\end{lemma}
\begin{proof}
\sketch For smooth $\Phi$, test the equation for $u$ with
$\Phi'(u)\varphi$. If $\Phi'\geq0$ and $\Phi''\geq0$, then
\[
\int a_{ij}D_i\Phi(u)D_j\varphi
=\int a_{ij}D_i uD_j(\Phi'(u)\varphi)
-\int a_{ij}D_i uD_j u\,\Phi''(u)\varphi\leq0.
\]
The supersolution case is analogous. Approximate a general convex $C^1$
function by convolution and pass to the limit by dominated convergence.
\end{proof}

\begin{remark}\label{hanlin:weaksol-part2-positive-part-subsolution}\dcref{Remark 4.7}\group{ch04-weak-solutions-part-ii}
\source{han-lin-lecture-notes:page-0084}
\uses{hanlin:weaksol-part2-convex-composition}
If $u$ is a subsolution, then $(u-k)^+=\max\{0,u-k\}$ is a subsolution.
\end{remark}

\begin{lemma}\label{hanlin:weaksol-part2-poincare-sobolev-zero-set}\dcref{Lemma 4.8}\group{ch04-weak-solutions-part-ii}
\source{han-lin-lecture-notes:page-0084,han-lin-lecture-notes:page-0085}
For $\epsilon>0$, if $u\in H^1(B_1)$ and
\[
|\{x\in B_1:u=0\}|\geq\epsilon|B_1|,
\]
then
\[
\int_{B_1}u^2\leq C\int_{B_1}|Du|^2,
\]
where $C$ depends only on $\epsilon$ and $n$.
\end{lemma}
\begin{proof}
\sketch Otherwise there are $u_m\in H^1(B_1)$ with zero sets of measure
at least $\epsilon|B_1|$, $\|u_m\|_{L^2}=1$, and
$\|Du_m\|_{L^2}\to0$. Compactness gives strong $L^2$ convergence to a
nonzero constant $u_0$, but
\[
\int_{B_1}|u_m-u_0|^2
\geq|u_0|^2|\{u_m=0\}|\geq\epsilon|u_0|^2|B_1|,
\]
a contradiction.
\end{proof}

\begin{theorem}\label{hanlin:weaksol-part2-density-theorem}\dcref{Theorem 4.9}\group{ch04-weak-solutions-part-ii}
\source{han-lin-lecture-notes:page-0085,han-lin-lecture-notes:page-0086}
\uses{hanlin:weaksol-part2-local-boundedness, hanlin:weaksol-part2-convex-composition, hanlin:weaksol-part2-poincare-sobolev-zero-set}
Let $\epsilon\in(0,1)$ and let $u$ be a positive supersolution in $B_2$.
If
\[
|\{x\in B_1:u\geq1\}|\geq\epsilon|B_1|,
\]
then
\[
\inf_{B_{1/2}}u\geq C,
\]
where $C>0$ depends only on $\epsilon,n$, and $\Lambda/\lambda$.
\end{theorem}
\begin{proof}
\sketch First replace $u$ by $u+\delta$ and later let $\delta\downarrow0$.
Lemma~\ref{hanlin:weaksol-part2-convex-composition} makes
$v=(\log u)^-$ a subsolution, so
Theorem~\ref{hanlin:weaksol-part2-local-boundedness} and
Lemma~\ref{hanlin:weaksol-part2-poincare-sobolev-zero-set} give
\begin{equation}
\sup_{B_{1/2}}v
\leq C\left(\int_{B_1}|Dv|^2\right)^{1/2}.
\end{equation}
Testing the supersolution inequality with $\zeta^2/u$ yields
\[
\int\zeta^2|D\log u|^2\leq C\int|D\zeta|^2.
\]
Choose $\zeta=1$ on $B_1$ and supported in $B_2$. The local boundedness
estimate then controls
$\sup_{B_{1/2}}(\log u)^-$ uniformly, giving
$\inf_{B_{1/2}}u\geq e^{-C}$.
\end{proof}

\begin{theorem}\label{hanlin:weaksol-part2-oscillation-decay}\dcref{Theorem 4.10}\group{ch04-weak-solutions-part-ii}
\source{han-lin-lecture-notes:page-0086}
\uses{hanlin:weaksol-part2-density-theorem}
If $u$ is a bounded solution of $Lu=0$ in $B_2$, then
\[
\osc_{B_{1/2}}u\leq\gamma\osc_{B_1}u,
\]
where $\gamma\in(0,1)$ depends only on $n$ and $\Lambda/\lambda$.
\end{theorem}
\begin{proof}
\sketch Let $\alpha_1=\sup_{B_1}u$ and $\beta_1=\inf_{B_1}u$. At least
one of the sets
\[
\left\{u\geq\frac{\alpha_1+\beta_1}{2}\right\},
\qquad
\left\{u\leq\frac{\alpha_1+\beta_1}{2}\right\}
\]
has measure at least $|B_1|/2$. Apply
Theorem~\ref{hanlin:weaksol-part2-density-theorem} to the corresponding
normalized nonnegative solution
\[
\frac{2(u-\beta_1)}{\alpha_1-\beta_1}
\quad\text{or}\quad
\frac{2(\alpha_1-u)}{\alpha_1-\beta_1}.
\]
In the first case the infimum on $B_{1/2}$ rises by a fixed fraction of
$\alpha_1-\beta_1$; in the second the supremum falls by that fraction.
Thus the oscillation contracts by $\gamma=1-1/C$.
\end{proof}

\begin{theorem}\label{hanlin:weaksol-part2-degiorgi-holder}\dcref{Theorem 4.11}\group{ch04-weak-solutions-part-ii}
\source{han-lin-lecture-notes:page-0087}
\uses{hanlin:weaksol-part2-local-boundedness, hanlin:weaksol-part2-oscillation-decay}
If $u\in H^1(B_1)$ solves $Lu=0$ in $B_1$, then
$u\in C^\alpha(B_1)$ for some $\alpha\in(0,1)$ depending only on $n$ and
$\Lambda/\lambda$. Moreover,
\[
\sup_{B_{1/2}}|u(x)|+
\sup_{x,y\in B_{1/2}}\frac{|u(x)-u(y)|}{|x-y|^\alpha}
\leq C\|u\|_{L^2(B_1)},
\]
where $C$ depends only on $n$ and $\Lambda/\lambda$.
\end{theorem}

\begin{lemma}\label{hanlin:weaksol-part2-energy-decay}\dcref{Lemma 4.12}\group{ch04-weak-solutions-part-ii}
\source{han-lin-lecture-notes:page-0087,han-lin-lecture-notes:page-0088}
\uses{hanlin:weaksol-part2-degiorgi-holder}
Suppose $a_{ij}\in L^\infty(B_r)$ satisfies
\[
\lambda|\xi|^2\leq a_{ij}(x)\xi_i\xi_j\leq\Lambda|\xi|^2
\qquad(x\in B_r,\ \xi\in\R^n),
\]
and $u\in H^1(B_r)$ satisfies
\[
\int_{B_r}a_{ij}D_i uD_j\varphi=0
\qquad(\varphi\in H^1_0(B_r)).
\]
Then there is $\alpha\in(0,1)$ such that, for $0<\rho<r$,
\[
\int_{B_\rho}|Du|^2
\leq C\left(\frac\rho r\right)^{n-2+2\alpha}
\int_{B_r}|Du|^2,
\]
where $C$ and $\alpha$ depend only on $n$ and $\Lambda/\lambda$.
\end{lemma}
\begin{proof}
\sketch Scale to $r=1$, subtract the mean of $u$, and use Poincar\'e's
inequality. Theorem~\ref{hanlin:weaksol-part2-degiorgi-holder} gives
\[
|u(x)-u(0)|^2\leq C|x|^{2\alpha}\int_{B_1}|Du|^2
\qquad(|x|\leq1/2).
\]
For $\rho\leq1/4$, test with
$\zeta^2(u-u(0))$, where $\zeta=1$ on $B_\rho$ and is supported in
$B_{2\rho}$. Caccioppoli's inequality gives
\[
\int_{B_\rho}|Du|^2
\leq C\rho^{n-2}\sup_{B_{2\rho}}|u-u(0)|^2,
\]
which proves the claim; larger $\rho$ are immediate.
\end{proof}

\begin{theorem}\label{hanlin:weaksol-part2-general-linear-holder}\dcref{Theorem 4.13}\group{ch04-weak-solutions-part-ii}
\source{han-lin-lecture-notes:page-0088}
\uses{hanlin:weaksol-part1-holder-sol, hanlin:weaksol-part2-energy-decay}
Assume $a_{ij}\in L^\infty(B_1)$ and $c\in L^n(B_1)$ satisfy
\[
\lambda|\xi|^2\leq a_{ij}(x)\xi_i\xi_j\leq\Lambda|\xi|^2
\qquad(x\in B_1,\ \xi\in\R^n).
\]
Suppose $u\in H^1(B_1)$ satisfies
\[
\int_{B_1}a_{ij}D_j uD_i\varphi+cu\varphi
=\int_{B_1}f\varphi
\qquad(\varphi\in H^1_0(B_1)).
\]
If $f\in L^q(B_1)$ for some $q>n/2$, then
$u\in C^\alpha(B_1)$ for some $\alpha\in(0,1)$ depending only on
$n,q,\lambda$, and $\Lambda$. Moreover, there is $R_0\in(0,1)$ such
that, for $x\in B_{1/2}$ and $r\leq R_0$,
\[
\int_{B_r(x)}|Du|^2
\leq Cr^{n-2+2\alpha}
\left(\|f\|_{L^q(B_1)}^2+\|u\|_{H^1(B_1)}^2\right),
\]
where $R_0$ and $C$ depend only on
$n,q,\lambda,\Lambda$, and $\|c\|_{L^n}$.
\end{theorem}

\section{Moser's Harnack Inequality}

Let $\Omega\subset\R^n$ be a domain and suppose
\[
\lambda|\xi|^2\leq a_{ij}(x)\xi_i\xi_j\leq\Lambda|\xi|^2
\qquad(x\in\Omega,\ \xi\in\R^n).
\]

\begin{theorem}\label{hanlin:weaksol-part2-scaled-local-boundedness}\dcref{Theorem 4.14}\group{ch04-weak-solutions-part-ii}
\source{han-lin-lecture-notes:page-0088}
\uses{hanlin:weaksol-part2-local-boundedness}
Let $u\in H^1(\Omega)$ be a nonnegative subsolution:
\[
\int_\Omega a_{ij}D_i uD_j\varphi
\leq\int_\Omega f\varphi
\]
for every nonnegative $\varphi\in H^1_0(\Omega)$. If
$f\in L^q(\Omega)$ for some $q>n/2$, then, for every
$B_R\subset\Omega$, $0<r<R$, and $p>0$,
\[
\sup_{B_r}u
\leq C\left((R-r)^{-n/p}\|u^+\|_{L^p(B_R)}
+R^{2-n/q}\|f\|_{L^q(B_R)}\right),
\]
where $C$ depends only on $n,\lambda,\Lambda,p$, and $q$.
\end{theorem}

\begin{theorem}\label{hanlin:weaksol-part2-weak-harnack}\dcref{Theorem 4.15}\group{ch04-weak-solutions-part-ii}
\source{han-lin-lecture-notes:page-0088,han-lin-lecture-notes:page-0089,han-lin-lecture-notes:page-0090,han-lin-lecture-notes:page-0091,han-lin-lecture-notes:page-0092,han-lin-lecture-notes:page-0093,han-lin-lecture-notes:page-0094}
\uses{hanlin:weaksol-part1-john-nirenberg, hanlin:weaksol-part2-local-boundedness}
Let $u\in H^1(\Omega)$ be a nonnegative supersolution:
\[
(*)\qquad
\int_\Omega a_{ij}D_i uD_j\varphi
\geq\int_\Omega f\varphi
\]
for every nonnegative $\varphi\in H^1_0(\Omega)$. If
$f\in L^q(\Omega)$ for some $q>n/2$, then, for every
$B_R\subset\Omega$, $0<p<n/(n-2)$, and $0<\theta<\tau<1$,
\[
\inf_{B_{\theta R}}u+R^{2-n/q}\|f\|_{L^q(B_R)}
\geq C\left(\frac1{R^n}\int_{B_{\tau R}}u^p\right)^{1/p},
\]
where $C$ depends only on $n,p,q,\lambda,\Lambda,\theta$, and $\tau$.
\end{theorem}
\begin{proof}
\sketch Scale to $R=1$. Put $\bar u=u+k>0$, where
$k=\|f\|_{L^q}$ if $f\not\equiv0$, and set $v=\bar u^{-1}$. Testing $(*)$
with $\bar u^{-2}\varphi$ shows that $v$ is a nonnegative subsolution with
a lower-order coefficient bounded in $L^q$.
Theorem~\ref{hanlin:weaksol-part2-local-boundedness} therefore gives
\begin{equation}
\inf_{B_\theta}\bar u
\geq C\left(\int_{B_\tau}\bar u^{-p}\right)^{-1/p}.
\end{equation}
It remains to control positive and negative moments of $\bar u$.

Let $w=\log\bar u$. Testing $(*)$ with $\bar u^{-1}\varphi^2$ gives
\begin{equation}
\int|Dw|^2\varphi^2\leq C\int|D\varphi|^2.
\end{equation}
For $B_{2r}(y)\subset B_1$, choose a cutoff supported in $B_{2r}(y)$ and
equal to one on $B_r(y)$. Then
\[
\int_{B_r(y)}|Dw|^2\leq Cr^{n-2},
\]
so Poincar\'e's inequality gives
\[
\frac1{r^n}\int_{B_r(y)}|w-w_{y,r}|\leq C.
\]
Thus $w\in\BMO$, and
Theorem~\ref{hanlin:weaksol-part1-john-nirenberg} supplies $p_0>0$ such that
\begin{equation}
\int_{B_\tau}e^{p_0|w-w_{B_\tau}|}\leq C,
\qquad
\left(\int_{B_\tau}\bar u^{-p_0}\right)
\left(\int_{B_\tau}\bar u^{p_0}\right)\leq C.
\end{equation}
Combining the local boundedness estimate with the moment product proves the
claim for $p_0$.

For $0<\gamma<1$, test $(*)$ with $\bar u^{\gamma-1}\eta^2$. With
$z=\bar u^{\gamma/2}$ and $\chi=n/(n-2)$, the energy and Sobolev
inequalities give, for $0<r<R<1$,
\begin{equation}
\left(\int_{B_r}\bar u^{\gamma\chi}\right)^{1/(\gamma\chi)}
\leq
\left(\frac{C}{(1-\gamma)^a(R-r)^2}\right)^{1/\gamma}
\left(\int_{B_R}\bar u^\gamma\right)^{1/\gamma}.
\end{equation}
Finitely iterating this reverse H\"older estimate upgrades $p_0$ to every
$p<n/(n-2)$. Undoing the scale and letting $k\downarrow0$ when
$f\equiv0$ proves the theorem.
\end{proof}

\begin{theorem}\label{hanlin:weaksol-part2-moser-harnack}\dcref{Theorem 4.16}\group{ch04-weak-solutions-part-ii}
\source{han-lin-lecture-notes:page-0094,han-lin-lecture-notes:page-0095}
\uses{hanlin:weaksol-part2-scaled-local-boundedness, hanlin:weaksol-part2-weak-harnack}
Let $u\in H^1(\Omega)$ be a nonnegative solution of
\[
\int_\Omega a_{ij}D_i uD_j\varphi
=\int_\Omega f\varphi
\qquad(\varphi\in H^1_0(\Omega)).
\]
If $f\in L^q(\Omega)$ for some $q>n/2$, then, for every
$B_R\subset\Omega$,
\[
\sup_{B_{R/2}}u
\leq C\left(\inf_{B_{R/2}}u
+R^{2-n/q}\|f\|_{L^q(B_R)}\right),
\]
where $C$ depends only on $n,\lambda,\Lambda$, and $q$.
\end{theorem}

\begin{corollary}\label{hanlin:weaksol-part2-moser-holder}\dcref{Corollary 4.17}\group{ch04-weak-solutions-part-ii}
\source{han-lin-lecture-notes:page-0095,han-lin-lecture-notes:page-0096}
\uses{hanlin:weaksol-part2-scaled-local-boundedness, hanlin:weaksol-part2-moser-harnack, hanlin:weaksol-part2-oscillation-iteration}
Let $u\in H^1(\Omega)$ solve
\[
\int_\Omega a_{ij}D_i uD_j\varphi
=\int_\Omega f\varphi
\qquad(\varphi\in H^1_0(\Omega)).
\]
If $f\in L^q(\Omega)$ for some $q>n/2$, then
$u\in C^\alpha(\Omega)$ for some $\alpha\in(0,1)$ depending only on
$n,q,\lambda$, and $\Lambda$. Moreover, for every $B_R\subset\Omega$ and
$x,y\in B_{R/2}$,
\[
|u(x)-u(y)|
\leq C\left(\frac{|x-y|}{R}\right)^\alpha
\left\{
\left(\frac1{R^n}\int_{B_R}u^2\right)^{1/2}
+R^{2-n/q}\|f\|_{L^q(B_R)}\right\},
\]
where $C$ depends only on $n,\lambda,\Lambda$, and $q$.
\end{corollary}
\begin{proof}
\sketch Scale to $R=1$. Put
$M(r)=\sup_{B_r}u$, $m(r)=\inf_{B_r}u$, and
$\omega(r)=M(r)-m(r)$. Apply
Theorem~\ref{hanlin:weaksol-part2-moser-harnack} to
$M(r)-u$ and $u-m(r)$ in $B_r$. Adding the two resulting inequalities gives
\[
\omega(r/2)\leq\gamma\omega(r)
+Cr^{2-n/q}\|f\|_{L^q(B_r)}
\]
for some $\gamma<1$.
Lemma~\ref{hanlin:weaksol-part2-oscillation-iteration} yields
\[
\omega(\rho)\leq C\rho^\alpha
\left(\omega(1/2)+\|f\|_{L^q(B_1)}\right).
\]
Theorem~\ref{hanlin:weaksol-part2-scaled-local-boundedness} bounds
$\omega(1/2)\leq C(\|u\|_{L^2(B_1)}+\|f\|_{L^q(B_1)})$.
Rescaling proves the stated estimate.
\end{proof}

\begin{lemma}\label{hanlin:weaksol-part2-oscillation-iteration}\dcref{Lemma 4.18}\group{ch04-weak-solutions-part-ii}
\source{han-lin-lecture-notes:page-0096}
Let $\omega$ and $\sigma$ be nondecreasing on $(0,R]$. Suppose
$0<\gamma,\tau<1$ and
\[
\omega(\tau r)\leq\gamma\omega(r)+\sigma(r)
\qquad(r\leq R).
\]
Then, for $\mu\in(0,1)$ and $r\leq R$,
\[
\omega(r)\leq C\left\{
\left(\frac rR\right)^\alpha\omega(R)
+\sigma(r^\mu R^{1-\mu})\right\},
\]
where $C$ depends only on $\gamma,\tau$, and
$\alpha=(1-\mu)\log\gamma/\log\tau$.
\end{lemma}
\begin{proof}
\sketch Fix $r_1\leq R$. Monotonicity and iteration give
\[
\omega(\tau^kr_1)
\leq\gamma^k\omega(R)+\frac{\sigma(r_1)}{1-\gamma}.
\]
For $r\leq r_1$, choose $k$ with
$\tau^kr_1<r\leq\tau^{k-1}r_1$. Then
\[
\omega(r)\leq\frac1\gamma
\left(\frac r{r_1}\right)^{\log\gamma/\log\tau}\omega(R)
+\frac{\sigma(r_1)}{1-\gamma}.
\]
Set $r_1=r^\mu R^{1-\mu}$.
\end{proof}

\begin{corollary}\label{hanlin:weaksol-part2-liouville}\dcref{Corollary 4.19}\group{ch04-weak-solutions-part-ii}
\source{han-lin-lecture-notes:page-0096,han-lin-lecture-notes:page-0097}
\uses{hanlin:weaksol-part2-moser-holder}
Suppose $u$ solves a homogeneous equation on $\R^n$:
\[
\int_{\R^n}a_{ij}D_i uD_j\varphi=0
\qquad(\varphi\in H^1_0(\R^n)).
\]
If $u$ is bounded, then it is constant.
\end{corollary}
\begin{proof}
\sketch The oscillation argument for
Corollary~\ref{hanlin:weaksol-part2-moser-holder} gives
$\omega(r)\leq\gamma\omega(2r)$ for some $\gamma<1$. Hence
\[
\omega(r)\leq\gamma^k\omega(2^kr)\leq C\gamma^k.
\]
Letting $k\to\infty$ shows $\omega(r)=0$ for every $r>0$.
\end{proof}

\section{Nonlinear Equations}

If $v=\Phi(u)$ with inverse $u=\Psi(v)$, then the change of test function
$\eta=\Psi'(v)\xi$ transforms
\[
\int a_{ij}D_i uD_j\xi=\int f\xi
\]
into
\[
\int a_{ij}D_i vD_j\eta
=\int\left(
\frac{\Psi''(v)}{\Psi'(v)}a_{ij}D_i vD_j v
+\frac{f}{\Psi'(v)}\right)\eta.
\]
Thus the transformed lower-order term has quadratic growth in $Dv$.

\begin{example}\label{hanlin:weaksol-part2-nonlinear-counterexample}\dcref{Example 4.20}\group{ch04-weak-solutions-part-ii}
\source{han-lin-lecture-notes:page-0097}
For $R<1$, the equation
\[
-\Delta u=|Du|^2\qquad\text{in }B_R\subset\R^2
\]
has both the zero solution and the unbounded weak solution
\[
u(x)=\log\log|x|^{-1}-\log\log R^{-1}
\in H^1(B_R)
\]
with zero boundary data. Thus boundedness is essential for uniqueness in the
class considered below.
\end{example}

Assume
\[
\lambda|\xi|^2\leq a_{ij}(x)\xi_i\xi_j\leq\Lambda|\xi|^2
\qquad(x\in B_1,\ \xi\in\R^n),
\]
and let $u\in H^1(B_1)\cap L^\infty(B_1)$ solve
\[
(\ast)\qquad
\int a_{ij}(x)D_i uD_j\varphi
=\int b(x,u,Du)\varphi
\]
for every $\varphi\in H^1_0(B_1)\cap L^\infty(B_1)$. The natural growth
condition is
\[
|b(x,u,p)|\leq C(u)(f(x)+|p|^2)
\qquad((x,u,p)\in B_1\times\R\times\R^n),
\]
where $f\in L^q(B_1)$ for some $q\geq2n/(n+2)$.

\begin{lemma}\label{hanlin:weaksol-part2-nonlinear-harnack}\dcref{Lemma 4.21}\group{ch04-weak-solutions-part-ii}
\source{han-lin-lecture-notes:page-0098,han-lin-lecture-notes:page-0099}
\uses{hanlin:weaksol-part2-scaled-local-boundedness, hanlin:weaksol-part2-weak-harnack}
Suppose $u\in H^1(B_1)$ is a nonnegative solution of $(\ast)$,
$|u|\leq M$ in $B_1$, and $b$ satisfies the natural growth condition with
$f\in L^q(B_1)$ for some $q>n/2$. Then, for every $B_R\subset B_1$,
\[
\sup_{B_{R/2}}u
\leq C\left(\inf_{B_{R/2}}u
+R^{2-n/q}\|f\|_{L^q(B_R)}\right),
\]
where $C$ depends only on $n,\lambda,\Lambda,M$, and $q$.
\end{lemma}
\begin{proof}
\sketch For large $\alpha>0$, set
\[
v=\frac{e^{\alpha u}-1}{\alpha}.
\]
The natural growth bound and ellipticity imply, for every nonnegative test
function $\varphi$,
\begin{equation}
\int a_{ij}D_i vD_j\varphi\leq C\int f\varphi,
\end{equation}
because the negative term
$-\alpha a_{ij}D_i uD_j u$ absorbs the quadratic growth. Since $u$ and $v$
are quantitatively comparable on $[0,M]$,
Theorem~\ref{hanlin:weaksol-part2-scaled-local-boundedness} gives
\begin{equation}
\sup_{B_{R/2}}u
\leq C\left\{
\left(\frac1{R^n}\int_{B_R}u^p\right)^{1/p}
+R^{2-n/q}\|f\|_{L^q(B_R)}\right\}.
\end{equation}
Similarly, $w=(1-e^{-\alpha u})/\alpha$ is a supersolution.
Theorem~\ref{hanlin:weaksol-part2-weak-harnack}
therefore gives, for $0<p<n/(n-2)$,
\[
\left(\frac1{R^n}\int_{B_R}u^p\right)^{1/p}
\leq C\left(
\inf_{B_{R/2}}u+R^{2-n/q}\|f\|_{L^q(B_R)}\right).
\]
Combining this with the preceding subsolution estimate proves the claim.
\end{proof}

\begin{remark}\label{hanlin:weaksol-part2-nonlinear-gradient-estimate}\dcref{Remark 4.22}\group{ch04-weak-solutions-part-ii}
\source{han-lin-lecture-notes:page-0099}
\uses{hanlin:weaksol-part2-nonlinear-harnack}
Testing the energy inequality from the proof of
Lemma~\ref{hanlin:weaksol-part2-nonlinear-harnack} with
$\varphi=(u+M)\eta^2$ gives
\[
\int|Du|^2\eta^2
\leq C\int(|D\eta|^2+|f|\eta^2),
\]
where $C$ depends only on $n,\lambda,\Lambda$, and $M$. Hence $Du$ has an
interior $L^2$ bound in terms of these constants and
$\|f\|_{L^1(B_1)}$.
\end{remark}

\begin{corollary}\label{hanlin:weaksol-part2-nonlinear-holder}\dcref{Corollary 4.23}\group{ch04-weak-solutions-part-ii}
\source{han-lin-lecture-notes:page-0099}
\uses{hanlin:weaksol-part2-moser-holder, hanlin:weaksol-part2-nonlinear-harnack}
Suppose $u\in H^1(B_1)$ is a bounded solution of $(\ast)$ and $b$ satisfies
the natural growth condition with $f\in L^q(B_1)$ for some $q>n/2$.
Then $u\in C^\alpha_{\mathrm{loc}}(B_1)$ for some $\alpha\in(0,1)$
depending only on $n,\lambda,\Lambda,q$, and $\|u\|_{L^\infty}$. Moreover,
\[
|u(x)-u(y)|\leq C|x-y|^\alpha
\qquad(x,y\in B_{1/2}),
\]
where $C$ depends only on
$n,\lambda,\Lambda,q,\|u\|_{L^\infty(B_1)}$, and
$\|f\|_{L^q(B_1)}$.
\end{corollary}
\begin{proof}
\sketch Repeat the oscillation argument for
Corollary~\ref{hanlin:weaksol-part2-moser-holder}, replacing the linear
Harnack estimate by Lemma~\ref{hanlin:weaksol-part2-nonlinear-harnack}.
\end{proof}

\begin{theorem}\label{hanlin:weaksol-part2-nonlinear-gradient-holder}\dcref{Theorem 4.24}\group{ch04-weak-solutions-part-ii}
\source{han-lin-lecture-notes:page-0099,han-lin-lecture-notes:page-0100}
\uses{hanlin:weaksol-part1-campanato, hanlin:weaksol-part1-lemma-3-4, hanlin:weaksol-part1-harmonic-lemma, hanlin:weaksol-part1-grad-holder, hanlin:weaksol-part2-nonlinear-gradient-estimate, hanlin:weaksol-part2-nonlinear-holder}
Suppose $u\in H^1(B_1)$ is a bounded solution of $(\ast)$, $b$ satisfies
the natural growth condition with $f\in L^q(B_1)$ for some $q>n$, and
$a_{ij}\in C^\alpha(B_1)$ for $\alpha=1-n/q$. Then
$Du\in C^\alpha_{\mathrm{loc}}(B_1)$ and
\[
\|Du\|_{C^\alpha(B_{1/2})}\leq C,
\]
where $C$ depends only on
$n,\lambda,\Lambda,q,\|u\|_{L^\infty(B_1)}$, and
$\|f\|_{L^q(B_1)}$.
\end{theorem}
\begin{proof}
\sketch It is enough to prove $Du\in L^\infty_{\mathrm{loc}}$, after which
Theorem~\ref{hanlin:weaksol-part1-grad-holder} applies. For
$B_r(x_0)\subset B_1$, let
$w-u\in H^1_0(B_r(x_0))$ solve the frozen-coefficient homogeneous equation
and put $v=u-w$. The maximum principle and
Lemma~\ref{hanlin:weaksol-part1-harmonic-lemma} give
\begin{equation}
\sup_{B_r(x_0)}|u-w|\leq\osc_{B_r(x_0)}u,
\end{equation}
and, for $0<\rho\leq r$,
\begin{equation}
\int_{B_\rho(x_0)}|Du-(Du)_{x_0,\rho}|^2
\leq C\left\{
\left(\frac\rho r\right)^{n+2}
\int_{B_r(x_0)}|Du-(Du)_{x_0,r}|^2
+\int_{B_r(x_0)}|Dv|^2\right\}.
\end{equation}
Testing the equation for $v$ by $v$ and using the bound on $u-w$ yields
\begin{equation}
\int_{B_r(x_0)}|Dv|^2
\leq C\left\{
\bigl(r^{2\alpha}+\osc_{B_r(x_0)}u\bigr)
\int_{B_r(x_0)}|Du|^2
+r^{n+2\alpha}\|f\|_{L^q(B_1)}^2\right\}.
\end{equation}
Corollary~\ref{hanlin:weaksol-part2-nonlinear-holder} first gives
$u\in C^{\delta_0}$ for some $\delta_0>0$. Combining the error estimate
with the energy estimate for $w$ and applying
Lemma~\ref{hanlin:weaksol-part1-lemma-3-4} gives,
for every $\delta<1$,
\[
\int_{B_r(x_0)}|Du|^2
\leq Cr^{n-2+2\delta}
\left(\int_{B_{7/8}}|Du|^2+\|f\|_{L^q(B_1)}^2\right).
\]
Thus $u\in C^\delta_{\mathrm{loc}}$ for every $\delta<1$. Choose
$\delta$ with $3\delta>2$ and $\alpha+\delta>1$. The error estimate gives
\[
\int_{B_r(x_0)}|Dv|^2\leq Cr^{n+2\alpha'}
\]
for some $\alpha'<\alpha$. The comparison estimate,
Lemma~\ref{hanlin:weaksol-part1-lemma-3-4}, and
Theorem~\ref{hanlin:weaksol-part1-campanato} imply
$Du\in C^{\alpha'}_{\mathrm{loc}}$, hence
$Du\in L^\infty_{\mathrm{loc}}$.
Theorem~\ref{hanlin:weaksol-part1-grad-holder} upgrades the exponent to
$\alpha$ and gives the stated estimate.
\end{proof}
